\section{A Publicly Tally-Hiding Schulze Election}\label{sec:schulze}
The scheme rests on three ideas. First, a ballot is a \emph{position
vector}, and its pairwise comparison bits are \emph{anchored} to that
vector (Section~\ref{sec:ballot}). Validity is then a conjunction of
statements about one ballot, so the scheme needs no shuffle and no
transitivity check, and the pairwise counts are a homomorphic sum. Second, a
\emph{closure certificate} certifies the winner over committed counts
(Section~\ref{sec:tally}). It consists of $O(m^3)$ $\Sigma$-proofs that
characterise the min-max closure of the defeat matrix without decrypting
any margin. The public learns the winner and, from the witness paths the
certificate carries in clear, the direction of some head-to-head
defeats, but no margin. Third, two
\emph{vector arguments} over Pedersen commitments replace the $m + m^2$
disjunctive proofs of the first ballot (Section~\ref{sec:batched}). The
ballot's validity becomes one linear relation, one Hadamard product, and
bit-decompositions of committed vectors. A generic
$\Sigma$-protocol for the preimage of a linear map and one Bayer--Groth
zero argument prove them in $O(m^2 \log m)$ exponentiations instead of $O(m^3)$.
Theorem names refer to the Rocq development of Section~\ref{sec:mech}.

\subsection{Setting and threat model}\label{sec:threat}
An election has $m$ candidates, $N$ voters, a bulletin board, and a set of
trustees holding an ElGamal key $h = g^x$ over a group $\langle g\rangle$ of
prime order $q$. Voters cast ballots from their own devices. The adversary
reads everything on the board, may corrupt voters' devices (and so submit
arbitrary, in particular malformed, ballots), and may corrupt all trustees but
one. We claim:
\begin{itemize}
  \item \emph{universal verifiability}, unconditionally: every check that
  decides the result is keyless and recomputable from the board.
  \item \emph{ballot privacy}, under IND-CPA security of ElGamal (DDH) and
  at least one honest trustee.
  \item \emph{public tally-hiding with leakage} $\mathcal{L}$, in the
  sense of Kryvos~\cite{kryvos2022}: the public learns
  $\mathcal{L}(D) = (w, \Pi)$, the winner and the matrix $\Pi$ of witness
  paths of the certificate (Section~\ref{sec:tally}), both deterministic
  functions of the count matrix $D$. $\Pi$ reveals the direction of some
  head-to-head contests and no count (Section~\ref{sec:security}). The
  trustees learn $D$.
  \item \emph{receipt-freeness relative to} $\mathcal{L}$ against a
  coercer who sees the entire bulletin board: nothing published
  distinguishes a ballot that follows the coercer's instruction from one
  that does not, unless the change alters $\mathcal{L}(D)$.
\end{itemize}

\subsection{A run, end to end}\label{sec:run}
An election has five steps, each posted on the bulletin board.
\begin{enumerate}
  \item \emph{Setup.} The trustees publish the election key $h = g^x$
  (in a deployment, a product of per-trustee keys, each published with a
  Schnorr proof of knowledge; the mechanised development uses a single
  key, see Section~\ref{sec:discussion}).
  \item \emph{Voting.} Each voter's device encrypts the voter's ranking,
  candidate by candidate and pair by pair, and attaches non-interactive
  proofs that the ciphertexts are those of \emph{one} ranking
  (Sections~\ref{sec:ballot} and~\ref{sec:batched}).
  \item \emph{Verification.} Anyone checks every posted ballot with the
  public verifier. Ballots that fail are posted as rejected.
  \item \emph{Aggregation.} Anyone multiplies the accepted ballots'
  pairwise ciphertexts and obtains the encrypted count matrix
  $\mathrm{Enc}(D)$ (Section~\ref{sec:aggregation}).
  \item \emph{Tally.} The trustees decrypt $D$ among themselves, run the
  Schulze method in private, and post a \emph{certificate}: the winner
  $w$ and zero-knowledge proofs, over re-encryptions of $D$, that $w$ is
  the Schulze winner of the counts inside $\mathrm{Enc}(D)$
  (Section~\ref{sec:tally}). Anyone checks the certificate against the
  aggregate of step~4.
\end{enumerate}
At the end the board holds the ciphertexts, commitments, and proofs of
every ballot, the aggregate (which anyone can recompute), the
certificate, and the winner. No ranking and no margin ever appears in
clear. The witness paths are the only cleartext besides the winner.

\subsection{Ballot encoding}\label{sec:ballot}
The Schulze method depends on the ballots only through the count matrix
$D$, and $D$ is a sum over ballots of pairwise bits. A ballot can
therefore be cast as its pairwise bits, which aggregate homomorphically,
provided the voter proves that the bits come from one ranking. Our
ballot does this by anchoring the bits to an encrypted ranking.

A ranking with ties is a \rocq[\emph{position vector}]{encrypt_ballot} $s \in [0,m]^m$: $s_i$ is
the position of candidate $i$, smaller is better, equal positions are ties,
and $s_i = m$ means ``unranked'', as in Gentoo's and Debian's ballot formats. Any such
vector is a preorder, so no transitivity condition is needed. For every
ordered pair $(i,j)$ the ballot's contribution to the count $D_{ij}$ is
the bit
\[
  b_{ij} = [\,s_i < s_j\,] = [\,\delta_{ij} \in \{1,\ldots,m\}\,],
  \quad \delta_{ij} = s_j - s_i,
\]
where $[P]$ is $1$ if $P$ holds and $0$ otherwise. The voter publishes $\mathrm{Enc}(s_i)$ for every candidate with a
membership proof $s_i \in \{0,\ldots,m\}$ (an $(m{+}1)$-way disjunctive
encryption proof, \rocqfile[EncProof]{src/Crypto/EncProof.v} of Section~\ref{schn_proof}), and
$\mathrm{Enc}(b_{ij})$ for every ordered pair with a proof
(\rocq{verify_pair_cell}) of
\begin{equation}\label{eq:consistency}
  \begin{aligned}
    &(b_{ij} = 1 \wedge \delta_{ij} \in \{1,\ldots,m\})\\
    \vee\; &(b_{ij} = 0 \wedge \delta_{ij} \in \{-m,\ldots,0\}).
  \end{aligned}
\end{equation}
The \rocq[verifier]{verify_ballot} computes
$\mathrm{Enc}(\delta_{ij}) = \mathrm{Enc}(s_j)\cdot\mathrm{Enc}(s_i)^{-1}$
itself, and the voter never supplies it. If $\mathrm{Enc}(\delta_{ij})$
were an input of the ballot, a voter could choose $\delta_{AB} =
\delta_{BA} = 1$ and prove $b_{AB} = b_{BA} = 1$, a ballot that counts
twice. Because the bits are anchored to one position vector, they are the
bits of one ranking, as \rocq{ballot_sound} states
(Section~\ref{sec:mech}).

In the implementation, \eqref{eq:consistency} is a 0/1 proof on
$\mathrm{Enc}(b_{ij})$ together with a single membership proof on the
ciphertext $\mathrm{Enc}(b_{ij})^{K}\cdot\mathrm{Enc}(\delta_{ij})$ with
$K = 2m+1$ (\rocq{pair_shift}, \rocq{pair_combined}). Its plaintext $K b_{ij} + \delta_{ij}$ lies in
$\{-m,\ldots,0\}$ exactly when $b_{ij}=0$ and $\delta_{ij}\le 0$, and in
$\{K+1,\ldots,K+m\}$ exactly when $b_{ij}=1$ and $\delta_{ij}\ge 1$, since
the shift $K$ separates the two ranges. For $m = 3$ and $K = 7$, the
honest pairs $(b, \delta) = (1, 2)$ and $(0, -1)$ give $9 \in \{8, 9,
10\}$ and $-1 \in \{-3, \ldots, 0\}$, while the dishonest pairs $(1, 0)$,
$(1, -1)$, and $(0, 2)$ give $7$, $6$, and $2$, which lie in neither
range. A membership proof (\rocqfile[EncProof]{src/Crypto/EncProof.v},
Section~\ref{schn_proof}) shows that a ciphertext encrypts one of a
public list of values, as a disjunction of Chaum--Pedersen statements
with one branch per value. Its cost is linear in the length of the list,
so the $m^2$ proofs with $2m+1$ branches dominate the ballot. A \rocq[ballot]{ballot} therefore carries
$m + m^2$ ciphertexts (the implementation keeps the $m$ diagonal cells,
which always encrypt $0$), $m$ membership proofs with $m{+}1$ branches, and
$m^2$ cells with a two-branch 0/1 proof and a $(2m{+}1)$-branch membership
proof.

\begin{table}[t]
  \caption{The running example, ballot 1: $A > B = C$, i.e.\ $s = (0,1,1)$,
  $m = 3$. The last two rows are ballots the verifier rejects.}
  \label{tab:example-ballot}
  \centering\footnotesize
  \begin{tabular}{@{}lccl@{}}
    \toprule
    pair $(i,j)$ & $b_{ij}$ & $\delta_{ij} = s_j - s_i$ & disjunct of \eqref{eq:consistency} \\
    \midrule
    $(A,B)$ & 1 & $1$  & $\delta \in \{1,2,3\}$ \\
    $(A,C)$ & 1 & $1$  & $\delta \in \{1,2,3\}$ \\
    $(B,A)$ & 0 & $-1$ & $\delta \in \{-3,\ldots,0\}$ \\
    $(C,A)$ & 0 & $-1$ & $\delta \in \{-3,\ldots,0\}$ \\
    $(B,C)$ & 0 & $0$  & $\delta \in \{-3,\ldots,0\}$ \\
    $(C,B)$ & 0 & $0$  & $\delta \in \{-3,\ldots,0\}$ \\
    \midrule
    $(B,A)$ & 1 & $-1$ & rejected: $-1 \notin \{1,2,3\}$ \\
    $(B,C)$ & 1 & $0$  & rejected: $0 \notin \{1,2,3\}$ \\
    \bottomrule
  \end{tabular}
\end{table}

\textbf{Running example:} Table~\ref{tab:example-ballot} shows the first ballot of a three-candidate
election, $A > B = C$. The tie forces $b_{BC} = b_{CB} = 0$, because
claiming $b_{BC} = 1$ would need $\delta_{BC} = 0 \in \{1,2,3\}$.
Reversing a bit fails the same way: $b_{BA} = 1$ would need $-1 \in
\{1,2,3\}$. The artefact's example run includes both rejected ballots.

\subsection{Batching the ballot proofs}\label{sec:batched}
The ballot of Section~\ref{sec:ballot} spends almost all of its cost on
the $m^2$ membership proofs with $2m+1$ branches: about $2m^3$
disjunction branches, each a few exponentiations
(Section~\ref{sec:eval:ballot}). The disjunctions are needed because a
$\Sigma$-protocol proves one statement about one ciphertext at a time.
The library's vector arguments, mechanised for the Swiss Post shuffle
verifier (Section~\ref{schn_proof}), prove one statement about a whole
committed vector, and with the position-anchored encoding the ballot's
validity is a statement of this kind. We therefore give a second,
\emph{batched} ballot with the same ciphertexts and the same
aggregation, and with two arguments in place of the $m + m^2$
disjunctions.

The batched ballot restates validity with three derived vectors. Write $\delta_{ij} = s_j - s_i$ and $p_{ij} = b_{ij}\,\delta_{ij}$, and
let $v_{ij} = -\delta_{ij} + 2p_{ij} - b_{ij}$: if $b_{ij} = 0$ then
$v_{ij} = -\delta_{ij}$, and if $b_{ij} = 1$ then $v_{ij} = \delta_{ij} -
1$. So a bit vector $b$ is consistent with the positions $s$ exactly when
every $b_{ij}$ is a bit and every $v_{ij}$ is a non-negative integer below
$2^L$, for $L = \lceil \log_2(m+1) \rceil$. No upper bound on the
positions is needed, since the pairwise differences are then all small
and the positions relative to one candidate are small integers. Validity
is thus: $b \in \{0,1\}^{m^2}$, $p = b \circ \delta$ (an entry-wise
product), $\delta = Ms$ (linear, $M$ the $m^2 \times m$ difference
matrix whose row $(i, j)$ is $e_j - e_i$), and $v = \sum_{k<L} 2^k X^{(k)}$ for bit vectors $X^{(k)}
\in \{0,1\}^{m^2}$ (linear in $X$). Each condition is linear in the
committed vectors, a Hadamard product, or a bit-ness condition, and the
vector arguments prove conditions of these three kinds.

\textbf{Worked example:} Take the ballot of Table~\ref{tab:example-ballot}, $s = (0, 1, 1)$, so
$m = 3$ and $L = 2$ (values of $v$ below $4$). Cell $(A, B)$ has $\delta
= 1$, $b = 1$, $p = 1$, $v = -1 + 2 - 1 = 0$, bits $X = (0, 0)$. Cell
$(B, A)$ has $\delta = -1$, $b = 0$, $p = 0$, $v = 1$, $X = (1, 0)$. The
tie $(B, C)$ has $\delta = 0$, $b = 0$, $p = 0$, $v = 0$. In the two
rejected ballots of the table, claiming $b_{BA} = 1$ with $\delta_{BA} =
-1$ forces $p = -1$ and $v = 1 - 2 - 1 = -2$, and claiming $b_{BC} = 1$
with $\delta_{BC} = 0$ forces $v = -1$. Neither value is a sum of two
bits, so no $X$ satisfies the linear relation and the ballot fails. The
bit decomposition of $v$ replaces the $2m+1$ membership branches: a
cheating voter is caught by a linear relation with no solution instead
of by a disjunction with no true branch.

The batched ballot uses three tools. The first is the Pedersen vector
commitment. For a key $h_1, g_1, \ldots, g_n$ of
generators with unknown discrete logarithms, \rocq[$\mathrm{com}(v; \rho)$]{com} $=
h_1^{\rho} \prod_j g_j^{v_j}$ commits to a whole vector $v \in
\mathbb{Z}_q^n$ in one group element. It is perfectly hiding, binding
under the discrete-logarithm assumption on the key, and homomorphic, $\mathrm{com}(v; \rho)\,\mathrm{com}(v'; \rho') =
\mathrm{com}(v + v'; \rho + \rho')$, so a verifier can derive a
commitment to any public linear combination of committed vectors, for
instance to $v - \mathbf{1}$ from a commitment to $v$. The same
homomorphism absorbs a public linear map into the key. With \rocq[$g'_i$]{gd} $=
\prod_k g_k^{M_{ki}}$, a commitment to $s$ under $(h_1, \mathbf{g}')$ is
a commitment to $Ms$ under $(h_1, \mathbf{g})$ with the same randomness
(\rocq{com_gd}), so $\delta = Ms$ needs no proof.

The second is the Bayer--Groth \rocq[zero argument]{zero_argument}. Given commitments to vectors $a_1, \ldots,
a_r$ and $b_1, \ldots, b_r$, it proves $\sum_i
a_i \star_y b_i = 0$ for a challenge $y$, where $a \star_y b = \sum_j
y^{j+1} a_j b_j$. As a polynomial in $y$, the left side has one
coefficient $\sum_i a_{ij} b_{ij}$ per coordinate $j$. It is the zero
polynomial exactly when every coordinate-wise sum of products vanishes,
and a non-zero polynomial of degree $n$ has at most $n$ roots. So an
argument accepted at $n$ distinct values of $y$ (the mechanised
statement) proves all $n$ coordinate relations, and a single random $y$
(the deployed statement) proves them except with probability $n/q$. With
$a = (b, p)$ and $b = (\delta, -\mathbf{1})$ the coordinate relations
are $b_j \delta_j - p_j = 0$, the Hadamard product. With $a = X$ and $b
= X - \mathbf{1}$ they are $X_j (X_j - 1) = 0$, bit-ness. One argument
carries several such relations if their rows are weighted with the
powers of a second challenge $z$. The sum is then a polynomial in $z$
whose coefficients are the individual relations, and it vanishes at $L +
2$ values of $z$ only if each relation does. The homomorphism applies
the weights to the commitments, $C \mapsto C^{z^i}$.

The third is a random linear combination. The ciphertexts $c_{ij}$ are
ElGamal, not Pedersen, so the vector arguments cannot speak about them
directly. Instead the verifier picks coefficients $\beta$ after the
ballot is fixed, everyone computes $C_\beta = \prod c_{ij}^{\beta_{ij}}$,
which encrypts $\langle \beta, b^{\mathrm{true}} \rangle$ for the true
plaintexts, and the voter proves that $C_\beta$ encrypts $\langle \beta,
b \rangle$ for the committed $b$. If $b \ne b^{\mathrm{true}}$ in some
coordinate, the two inner products agree only for a $1/q$ fraction of
the $\beta$, so one accepted link ties the committed vector to the
plaintexts except with probability $1/q$. This is the batch-verification
idea of Bellare, Garay, and Rabin~\cite{bgr98}, which the
shuffle argument of Bayer and Groth~\cite{10.1007/978-3-642-29011-4_17}
also uses. Here it feeds $\Sigma$-protocol vector arguments. It is the only place
where the batched ballot's soundness is probabilistic and not exact.

The voter publishes the ciphertexts $c_{ij} = \mathrm{Enc}(b_{ij})$ and
$L + 3$ Pedersen vector commitments over one key $h_1, \mathbf{g}$ of $m^2$
generators: $C_\delta = \mathrm{com}_{\mathbf{g}'}(s; \rho_\delta)$ under
the derived key, $C_b$, $C_p$, and $C_{X^{(k)}}$ for $k < L$
(Figure~\ref{fig:batched}). The positions are committed rather than
encrypted, so the verifier no longer derives $\mathrm{Enc}(\delta_{ij})$,
and the ballot has $m^2$ ciphertexts instead of $m + m^2$. The verifier also forms (\rocq{C0})
\[
  C_0 = C_\delta\, C_b\, C_p^{-2} \prod_{k<L} C_{X^{(k)}}^{\,2^k},
\]
which by the homomorphism commits to $\delta + b - 2p + \sum_k 2^k
X^{(k)}$. This is the zero vector when $v = \sum_k 2^k X^{(k)}$, so the
voter proves that identity by opening $C_0$ as $h_1^{\rho^*}$ with
$\rho^* = \rho_\delta + \rho_b - 2\rho_p + \sum_k 2^k \rho_{X^{(k)}}$.
Two arguments remain.
\begin{enumerate}
  \item[(a)] \emph{Linear relations.} One $\Sigma$-protocol for the
  preimage of a linear map \rocq[$\varphi$]{phi3} $: \mathbb{Z}_q^W \to G^5$, $W = m +
  m^2 + 4$, applied to the witness $w = (s, b, \rho_\delta, \rho_b, R,
  \rho^*)$. Its five rows say that $C_\delta$ opens to $s$ under
  $\mathbf{g}'$, that $C_b$ opens to $b$, that the ciphertext
  combination $C_\beta = \prod_{ij} c_{ij}^{\beta_{ij}}$, for coefficients
  $\beta$ drawn by the verifier after the ballot is fixed, is the
  encryption of $\langle \beta, b \rangle$ with randomness $R$ (two rows,
  $g^R$ and $g^{\langle\beta,b\rangle} h^R$), and that $C_0 =
  h_1^{\rho^*}$. The response is one field element per witness
  coordinate.
  \item[(b)] \emph{One zero argument} with $L + 3$ rows weighted by the
  verifier's $z$: $z\,b$ against $\delta$, $z\,p$ against $-\mathbf{1}$,
  $z^{2+k} X^{(k)}$ against $X^{(k)} - \mathbf{1}$ for $k < L$, and
  $z^{2+L}\, b$ against $b - \mathbf{1}$, all commitments derived from the
  posted ones by the homomorphism. Its claim, $z\,(b \star_y \delta - p
  \star_y \mathbf{1}) + \sum_k z^{2+k}\, X^{(k)} \star_y (X^{(k)} -
  \mathbf{1}) + z^{2+L}\, b \star_y (b - \mathbf{1}) = 0$, forces $p = b
  \circ \delta$ and the bit-ness of every $X^{(k)}$ and of $b$ once it
  holds at $L + 2$ values of $z$ and $m^2$ values of $y$.
\end{enumerate}
The zero argument is Bayer and Groth's (the sub-argument of the Hadamard
argument of their shuffle), with the bilinear map $\star_y$ of the Swiss
Post specification. The library has its \rocq[completeness]{zero_completeness}, its
\rocq[$(2r{+}1)$-special soundness]{zero_soundness} for $r$ rows, and its \rocq[perfect
honest-verifier zero-knowledge]{zero_shvzk} (Section~\ref{schn_proof}).

\begin{figure}[t]
  \centering\footnotesize
  \fbox{\begin{minipage}{0.96\columnwidth}
  \textbf{Public:} $g, h$ (election key), $h_1, \mathbf{g} \in G^{m^2}$
  (commitment key), $m$, $L$; derived $g'_i = \prod_k g_k^{M_{ki}}$.\\
  \textbf{Voter's witness:} positions $s \in \mathbb{N}^m$; derived
  $\delta = Ms$, $b = [\delta > 0]$, $p = b \circ \delta$, $X^{(k)}$ the bits of
  $v = -\delta + 2p - b$; randomness $r \in \mathbb{Z}_q^{m^2}$ for the
  ciphertexts, $\rho_\delta, \rho_b, \rho_p, \rho_{X^{(k)}}$ for the
  commitments.\\[2pt]
  \textbf{Message 1 (the ballot):} $c_{ij} = \mathrm{Enc}(b_{ij}; r_{ij})$;
  $C_\delta = \mathrm{com}_{\mathbf{g}'}(s; \rho_\delta)$, $C_b$, $C_p$,
  $C_{X^{(k)}}$ ($k < L$) over $\mathbf{g}$.\\
  \textbf{Challenge 1:} $\beta \in \mathbb{Z}_q^{m^2}$; $y, z \in \mathbb{Z}_q$.\\
  \textbf{Message 2 (announcements):} $A = \varphi_{\beta}(a)$ for
  random $a \in \mathbb{Z}_q^W$; the announcement of the zero argument
  for $y, z$.\\
  \textbf{Challenge 2:} $c$ (linear proof), $x$ (zero argument).\\
  \textbf{Message 3 (responses):} $u = a + c\,w \in \mathbb{Z}_q^W$; the
  response of the zero argument.\\[2pt]
  \textbf{Verifier:} $|c| = m^2$; $C_\beta = \prod c_{ij}^{\beta_{ij}}$ and
  $C_0 = C_\delta C_b C_p^{-2} \prod_k C_{X^{(k)}}^{2^k}$;
  $\varphi_{\beta}(u) = A + c\cdot T$ where $T = (C_\delta, C_b, C_\beta,
  C_0)$, i.e.\ the five group equations
  $\mathrm{com}_{\mathbf{g}'}(u_s; u_{\rho_\delta}) = A_1 C_\delta^{c}$,
  $\mathrm{com}(u_b; u_{\rho_b}) = A_2 C_b^{c}$,
  $g^{u_R} = A_3\, C_{\beta,1}^{c}$,
  $g^{\langle\beta, u_b\rangle} h^{u_R} = A_4\, C_{\beta,2}^{c}$,
  $h_1^{u_{\rho^*}} = A_5\, C_0^{c}$; the zero argument at $(y, x)$ on the
  rows $(C_b^{z}, C_p^{z}, C_{X^{(k)}}^{z^{2+k}}, C_b^{z^{2+L}})$ against
  $(C_\delta,\ \mathrm{com}(-\mathbf{1}; 0),\ C_{X^{(k)}}
  \mathrm{com}(-\mathbf{1}; 0),\ C_b\, \mathrm{com}(-\mathbf{1}; 0))$.\\[2pt]
  \textbf{Fiat--Shamir:} $\beta, y, z$ are hashes of the ciphertexts and
  commitments (tagged, one query per coordinate of $\beta$); $c$ and $x$
  are hashes of the same and of the respective announcement. The honest
  voter computes the ballot in three passes (\rocq{nizk_encrypt_ballot3}).
  \end{minipage}}
  \caption{The batched ballot (\rocq{ballot3}, verifier \rocq{verify_ballot3}). $\mathrm{com}(v; \rho) = h_1^{\rho}
  \prod_j g_j^{v_j}$; $u_s, u_b, \ldots$ are the slices of $u$
  matching the witness layout.}
  \label{fig:batched}
\end{figure}

Argument (a) is an instance of a $\Sigma$-protocol for the preimage of a
linear map, which we add to the library
(\rocqfile[Crypto/LinearSigma.v]{src/Crypto/LinearSigma.v}): for a vector space $V$ over
$\mathbb{Z}_q$ and a linear map $\varphi : \mathbb{Z}_q^W \to V$, prove
knowledge of $w$ with $\varphi(w) = T$ by announcing $A = \varphi(a)$ for
a random $a$, receiving $c$, and answering $u = a + c\,w$. The verifier
checks $\varphi(u) = A + c\,T$ (\rocq{lin_verify}, completeness \rocq{lin_completeness}). This is Schnorr's protocol with the
exponentiation replaced by any linear map, and Okamoto's representation
proof is the case $V = G$. Special soundness is exact, $w = (u - u')/(c -
c')$ (\rocq{lin_soundness}), and honest-verifier zero-knowledge is
perfect, by the bijection $a \mapsto u$ on the randomness
(\rocq{lin_shvzk}, in the permutation form of
Section~\ref{schn_proof}). We need the generality because $V$ may be
a power or a product. Here $V = G^5$, with the library's \rocq[power vector
spaces]{vec_vspace}, so the two commitment openings, the ciphertext link, and
the opening of $C_0$ are one proof with one challenge and a response of
$W$ field elements, instead of five proofs tied together by hand.

The soundness theorem of the batched ballot has the perfect-adversary
form of the library's shuffle theorems. It assumes a prover that answers
every challenge the extractor needs, and the forking argument that turns
this into a probability is on paper.
\begin{quote}
  \emph{Theorem (\rocq{ballot3_sound}).} Let $h = g^x$, $L \ge 1$, and
  let the field have characteristic above $3 \cdot 2^L + 2$. Suppose a
  ballot is accepted (a) on a base vector $\beta$ with
  two challenges $c \ne c'$ on the same announcement, and, for every
  coordinate $k$, on a vector that differs from $\beta$ in coordinate $k$
  only, again with two challenges; and (b) on $m^2$ distinct non-zero
  values of $y$ and, for each, $L + 2$ distinct non-zero values of $z$,
  with $2(L+3)+1$ challenges $x$ each on one announcement. Then either
  the binding of $(h_1, \mathbf{g})$ is broken, or there is a position
  vector $\mathit{pos} \in [0, 2^{L+1}]^m$ such that every ciphertext
  $c_{ij}$ decrypts to $g^{[\mathit{pos}_i < \mathit{pos}_j]}$.
\end{quote}
The proof runs the library's extractors in sequence
(Appendix~\ref{app:batched}). The linear proof opens the commitments and
links $C_\beta$ to the committed $b$. Roots in the group turn the link
into $\mathrm{Dec}(c_k) = g^{b_k}$. The zero argument, at $L + 2$
values of $z$ and $m^2$ values of $y$, forces $p = b \circ \delta$ and
the bit-ness of $b$ and of every $X^{(k)}$, from which
\rocq{decode_positions} reads off the positions. The ciphertext link is
undone by roots and not by comparing exponents, so the batched ballot
needs no hypothesis on the order of $g$, while the first ballot needs
\rocq{exp_injective}.

The batched ballot is also perfectly zero-knowledge.
\begin{quote}
  \emph{Theorem (\rocq{ballot3_zk}).} For positions in $[0,m]$ and
  $m < 2^L$, the distribution of the honest ballot over the voter's
  randomness (the $m^2$ encryption randomnesses, the $L + 3$ commitment
  randomnesses, the announcement of the linear proof, and the randomness
  of the zero argument) is a permutation of the output of a
  simulator that has the ciphertexts and the challenges only: it draws
  the $L + 3$ commitments uniformly in $\langle h_1\rangle$ and runs the
  library's simulators of the linear proof and of the zero argument on
  them.
\end{quote}
The proof composes three bijections on the randomness
(Appendix~\ref{app:batched}). Like the library's Bayer--Groth simulators,
it names the discrete logarithms of the commitment key in base $h_1$,
which exist for every key in the subgroup and which the simulator does
not use. The
proofs and the commitments therefore reveal nothing about the ranking beyond what
the $m^2$ ciphertexts do, and receipt-freeness
(Section~\ref{sec:security}) carries over unchanged.

Made non-interactive with Fiat--Shamir, the ballot takes five oracle
queries where the first ballot takes one (Figure~\ref{fig:batched}). The
order matters for soundness. The challenges $\beta$, $y$, and $z$ must be
fixed after the ciphertexts and the commitments, since a voter who chose
$\beta$ could answer the link with a $b$ that is not the plaintext
vector, and $c$ and $x$ must be fixed after the respective announcements.
Attema, Fehr, and Klo{\ss}~\cite{afk21} bound the soundness loss of
Fiat--Shamir for such multi-round, $k$-special-sound protocols. The honest
voter's announcements do not depend on $c$ or $x$, so the derived
challenges are stable across the voter's three passes and the honest
ballot is accepted (\rocq{nizk_ballot3_valid}). Each coordinate of
$\beta$ is one tagged query, so the challenge vector is as long as the
ballot needs without any assumption on the output length of the oracle.

The prover computes the commitments and the two announcements, about
$(L + 4)m^2$ exponentiations in all, and the verifier about $5m^2$. The
ballot is $2m^2 + 3L + 17$ group elements, of which $2m^2$ are the
ciphertexts, and $3m^2 + m + 7$ field elements. At $m = 17$ the prover
takes $0.7$\,s, the verifier $1.6$\,s, and the ballot is $180$\,KB,
against $130$\,s, $68$\,s, and $6.2$\,MB for the first ballot
(Table~\ref{tab:ballotcost}). The ciphertexts are now $148$ of the
$180$\,KB. Table~\ref{tab:twoballots} compares the two ballots.
Section~\ref{sec:discussion} lists the costs of batching: soundness in
the rewinding form becomes statistical and multi-round, binding of the
commitment key is assumed, zero-knowledge names the discrete logarithms
of the key, and the ballot takes five oracle queries instead of one.
The development and the paper keep both ballots. The first has exact
special soundness from two transcripts and needs no key beyond the
election key. The second is the one to deploy. Both plug into the same
aggregation and the same certificate, with end-to-end theorems
\rocq{schulze_end_to_end} and \rocq{schulze_end_to_end3}.

\begin{table}[t]
  \caption{The two ballots for $m$ candidates, $L = \lceil \log_2 (m+1)
  \rceil$. Measured values are in Table~\ref{tab:ballotcost}.}
  \label{tab:twoballots}
  \centering\scriptsize\setlength{\tabcolsep}{3pt}
  \begin{tabularx}{\columnwidth}{@{}l>{\raggedright\arraybackslash}X>{\raggedright\arraybackslash}X@{}}
    \toprule
     & first ballot (\S\ref{sec:ballot}) & batched ballot (\S\ref{sec:batched}) \\
    \midrule
    ciphertexts & $m + m^2$ & $m^2$ \\
    commitments & --- & $L + 3$, one key of $m^2$ generators \\
    proofs & $m$ membership, $m^2$ 0/1, $m^2$ membership & $1$ linear, $1$ zero argument of $L + 3$ rows \\
    size & $O(m^3)$ group elements & $2m^2 + O(\log m)$ group, $3m^2 + O(m)$ field elements \\
    prover, verifier & $O(m^3)$ exp. & $O(m^2 \log m)$, $O(m^2)$ exp. \\
    oracle queries & $1$ & $5$ \\
    soundness & exact, $2$ transcripts & statistical, $O(m^2 \log^2 m)$ transcripts \\
    group hyp. & $g$ of order $q$, char.\ $> 4m+2$ & char.\ $> 3 \cdot 2^L + 2$ \\
    key hyp. & none & binding of the commitment key \\
    zero-knowledge & perfect, no hypothesis & perfect, key in $\langle h_1 \rangle$ \\
    Rocq & \rocq{ballot_sound}, \rocq{ballot_zk} & \rocq{ballot3_sound}, \rocq{ballot3_zk} \\
    \bottomrule
  \end{tabularx}
\end{table}

\subsection{Aggregation}\label{sec:aggregation}
The encrypted count matrix is the homomorphic product of the accepted
ballots' bits, $\mathrm{Enc}(D_{ij}) = \prod_{\text{ballots}}
\mathrm{Enc}(b_{ij})$, where $D_{ij}$ is the number of voters who rank $i$
strictly above $j$. Aggregation needs no key: anyone recomputes it from the
board (\rocq{aggregate_margins}, \rocq{aggregate_margins_enc}). The
counts are never decrypted on the board.

\subsection{Tally}\label{sec:tally}
Let $N$ be the number of accepted ballots, so $0 \le D_{ij} \le N$. The
trustees decrypt $D$ among themselves and compute the Schulze winner in
private. Instead of $D$ they publish a \rocq[\emph{certificate}]{tally_proof}: a proof,
checkable by anyone against the encrypted aggregate, that they ran the
Schulze method correctly. The certificate follows the method's three
steps (counts to defeats, defeats to strongest paths, strongest paths
to winner). For each step it publishes encryptions of the intermediate
values with proofs that they stand in the right relation to the
previous ones. Every relation of the Schulze method compares two
integers in $[0, N]$, and a comparison of encrypted integers is a
membership proof on the difference ciphertext: $\mathrm{Enc}(u) /
\mathrm{Enc}(u')$ encrypts $u - u'$, and $u \ge u'$ is $u - u' \in \{0,
\ldots, N\}$ (\rocq{verify_range}). The certificate needs no other
primitive, and this is the same \rocqfile[EncProof]{src/Crypto/EncProof.v}
as in the ballot. The certificate has four parts. Each is a
$\Sigma$-protocol of the library, and all share one Fiat--Shamir
challenge (\rocq{nizk_generate_tally}).
\begin{enumerate}
  \item[(i)] \emph{Re-encryption.} For every pair a fresh encryption
  $\mathrm{Enc}'(D_{ij})$ with a decryption proof (Chaum--Pedersen) that
  $\mathrm{Enc}'(D_{ij})/\mathrm{Enc}(D_{ij})$ decrypts to $1$. $D$ is
  never on the board in clear, and the rest of the certificate is about
  these re-encryptions.
  \item[(ii)] \emph{Defeats.} The defeat strength of $i$ over $j$ is
  $d_{ij} = D_{ij}$ if $D_{ij} > D_{ji}$ and $0$ otherwise. The trustees
  publish $\mathrm{Enc}(d_{ij})$ with a range proof $d_{ij} \in
  \{0,\ldots,N\}$ and a disjunction of two membership proofs on
  differences: either $d_{ij} = 0$ and $D_{ij} - D_{ji} \le 0$, or
  $d_{ij} = D_{ij}$ and $D_{ij} - D_{ji} \ge 1$ (implemented with the same
  shift trick as \eqref{eq:consistency}, $K = 2N+1$; \rocq{defeat_statements}, \rocq{tally_shift}).
  \item[(iii)] \emph{Strongest paths.} The trustees publish
  $\mathrm{Enc}(p_{ij})$ for the strongest-path matrix $p$ and prove that
  it is the min-max closure of $d$: $p_{ij} - d_{ij} \in \{0,\ldots,N\}$
  for every pair; for every $k$, $p_{ij} \ge p_{ik}$ or $p_{ij} \ge
  p_{kj}$ (a $(2N{+}2)$-branch disjunction over the pair of differences, \rocq{upper_statements});
  and for every pair a witness path $i = u_0, u_1, \ldots, u_\ell = j$ with
  a proof $d_{u_t u_{t+1}} - p_{ij} \in \{0,\ldots,N\}$ for every edge (\rocq{verify_path}, \rocq{verify_edges}).
  \item[(iv)] \emph{Winner.} For the announced winner $w$, a proof
  $p_{wj} - p_{jw} \in \{0,\ldots,N\}$ for every $j \ne w$.
\end{enumerate}
The first two conditions of (iii) say that $p$ is closed and dominates
$d$, hence dominates the closure. The witness paths say that a path
attains every $p_{ij}$, hence the closure dominates $p$. Together they
characterise the closure, and (iv) then says that $w$ is a Schulze
winner:
\begin{quote}
  \emph{Lemma (certificate soundness).} If $d$, $p$, and the witness paths
  satisfy (iii) (\rocq{cert_ok}), then $p$ is the min-max closure of $d$, the closure is
  unique, and a $w$ satisfying (iv) is a Schulze winner of $d$
  (\rocq{cert_ok_is_closure}, \rocq{closure_unique},
  \rocq{cert_winner}; the closure is the Kleene closure of the min-max
  semiring~\cite{10.1007/978-3-319-66107-0_26}).
\end{quote}
The strongest-path matrix is the closure of $d$ under ``a path through
$k$ is at least as strong as the weaker of its two halves'', that is,
the least matrix $p \ge d$ with $p_{ij} \ge \min(p_{ik}, p_{kj})$ for
all $k$. The first two conditions of (iii) say that the published $p$
is such a matrix, hence at least the closure. Without the third, the
trustees could publish any larger matrix, for instance all $N$, and
make any candidate a winner. The witness paths rule this out: a path of
strength $\ge p_{ij}$ exists, so the closure is at least $p_{ij}$. The
trustees find the paths when they compute the closure
(\rocq{witness_path}) and publish them in clear. Each path costs one
range proof per edge. In
the running example the witness for $(C, B)$ is $C \to A \to B$ with
edge strengths $3$ and $3$, and the closure disjunction for $(C, B)$
through $k = A$ reads $p_{CB} \ge p_{CA}$ or $p_{CB} \ge p_{AB}$, i.e.\
$3 \ge 3$ or $3 \ge 3$.

The verifier checks the proofs of (i)--(iv) against the aggregated
$\mathrm{Enc}(D)$ and the published commitments
(\rocq{verify_tally}, one check per part in \rocq{tally_checks}). The
certificate has $O(m^3)$ proofs, mostly the $m^3$ closure disjunctions,
and their number does not depend on $N$. Each range proof is a
membership proof with $N{+}1$ branches and each closure disjunction has
$2N{+}2$ branches, so the certificate has $O(m^3 N)$ group elements.
Section~\ref{sec:eval} measures both.

Gentoo's council and Debian's committees fill several seats from one
ranked election by removing the winner and running the method again on
the remaining candidates. The certificate supports this without touching
the ballots. The aggregate $\mathrm{Enc}(D)$ restricted to the remaining
candidates is the aggregate for the smaller election, and the trustees
certify a winner of it. Seven seats are seven certificates on the same
encrypted counts, of $m, m-1, \ldots, m-6$ candidates.

\textbf{Running example, continued:} With ballots $2$--$5$ equal to $B > C > A$, $B > C > A$, $C > A > B$, and
$A > C > B$, the pairwise counts are $D_{AB} = 3$, $D_{BA} = 2$, $D_{AC} = 2$,
$D_{CA} = 3$, and $D_{BC} = D_{CB} = 2$. The defeats are $d_{AB} = 3$ and
$d_{CA} = 3$, all others $0$. The strongest paths are $p_{AB} = 3$,
$p_{CA} = 3$, and $p_{CB} = 3$ (via $C \to A \to B$), all others $0$.
$C$ is the winner, since $p_{CA} = 3 \ge p_{AC} = 0$ and $p_{CB} = 3 \ge
p_{BC} = 0$. The board shows the five ballots, the nine aggregated
ciphertexts, and the certificate: nine re-encryptions with decryption
proofs, nine encrypted defeats with range proofs and defeat disjunctions,
nine encrypted path strengths with range proofs, $27$ closure
disjunctions, the witness paths with their edge proofs, and two winner
inequalities. The artefact's example program prints all of it and also
runs the certificate for $A$, which the verifier rejects. Of the nine
witness paths, only $C \to A \to B$ has two edges. It tells the public
that $C$ defeats $A$ and $A$ defeats $B$, but not by how much. Nothing on
the board distinguishes ballot $1$ ($A > B = C$) from ballot $5$ ($A > C >
B$). A coercer who demanded $A > B = C$ sees two encrypted ballots with
simulatable proofs and a certificate about the aggregate, and the
count matrix that would tell the two apart is never published.

\subsection{Security argument}\label{sec:security}
The scheme is universally verifiable because every check is keyless:
ballot verification, aggregation, and certificate verification take only
the board, and a rejected ballot is posted as rejected, so the verdict is
auditable (\rocq{count_sound}, Section~\ref{sec:mech}).

For soundness, \rocq{ballot_sound} states that two accepting transcripts
of a ballot with the same announcements and different challenges yield
a position vector $s$ such that every bit encrypts $[s_i < s_j]$.
\rocq{ballot3_sound} gives the same conclusion, or a binding break, for
the batched ballot from the transcripts of Section~\ref{sec:batched}.
\rocq{verify_tally_sound} states that two accepting transcripts of a
certificate on counts $D \in [0,N]$ yield $d = \mathrm{defeat}(D)$, a
$p$ satisfying (iii), and a winner satisfying (iv), hence a Schulze
winner by the lemma. \rocq{schulze_end_to_end} combines the two:
accepted ballots and an accepted certificate on their aggregate name a
Schulze winner of the rankings the voters cast. \rocq{schulze_end_to_end3}
does the same with batched ballots. The statements are in the rewinding
form of special soundness, and the Fiat--Shamir step is on paper, as for
the rest of the library (Section~\ref{sec:related}).

For privacy, \rocq{ballot_zk}, \rocq{ballot3_zk}, and \rocq{tally_zk}
state that the distribution of every ballot transcript and of every
certificate transcript is a permutation of the distribution produced by
a simulator that has neither the ranking nor the counts. The ballot
simulator has the ciphertexts and the challenges. The certificate
simulator also has the fresh ciphertexts of parts (i)--(iii) and the
witness paths $\Pi$, which the certificate publishes in clear. Ballot
privacy and public tally-hiding then need the hiding of the ciphertexts,
i.e.\ IND-CPA security of ElGamal under DDH and at least one honest
trustee. This is the standard on-paper reduction, and it is not
mechanised. What remains public is $\mathcal{L}(D) = (w, \Pi)$, and
$\Pi$ is computed from $D$ by \rocq{witness_path}. A path with one edge
is published both when $p_{ij} = 0$ and when $d_{ij} \ge p_{ij} > 0$, so
it does not say whether $i$ defeats $j$. A path with two or more edges
arises only when $p_{ij} > 0$, and every edge $(u, v)$ on it has $d_{uv}
\ge p_{ij} > 0$: it reveals that $u$ defeats $v$, but not by how much.
On the thirteen Gentoo council elections the paths reveal the
head-to-head winner of $118$ of the $1050$ candidate pairs, $24$ of them
among the $66$ pairs with a margin of at most two, and no margin
(Table~\ref{tab:leak}).

Receipt-freeness against a board-reading coercer follows from the same
facts. The coercer's view is the board. For the coerced voter and for
everyone else it holds $m^2$ ciphertexts (and, in the batched ballot,
$L + 3$ commitments) with proofs. It also holds the aggregate, which the
ballots determine, and the certificate, which is simulatable from its
ciphertexts, $w$, and $\Pi$. If the coerced voter's ranking is replaced
by any other, the ciphertexts are indistinguishable under DDH, the
proofs are simulated identically, and the coercer can tell the two
boards apart only if $\mathcal{L}(D)$ changes. That is more than the
winner alone reveals: a revealed defeat with a margin of one flips if one
voter reverses that preference. To test a voter this way, the coercer
must know that the margin is that small, and neither the margins nor
the counts are published, unlike in the systems that publish the count matrix
(Section~\ref{sec:eval:leak}). Nor can the voter prove how they voted,
since any ranking is consistent with their board entry. The argument does
not cover a coercer who obtains the device's encryption randomness, which
opens the ciphertexts. Protecting against that coercer is coercion
resistance, which the scheme does not provide. The trustees learn $D$ and are trusted for privacy but
not for correctness: a trustee that publishes a wrong winner fails the
certificate.

\subsection{Cost and comparison}\label{sec:comparison}
Table~\ref{tab:comparison} compares the scheme with other verifiable Schulze and
Condorcet tallies. The MPC-based systems also hide the counts from the
talliers, at a cost of hours to days for $20$ candidates. Kryvos is
publicly tally-hiding without our scheme's path leakage, but it needs a
SNARK with a trusted setup and supports Condorcet winners rather than
Schulze. The
two systems that publish the counts or the ballots are cheaper but leak
(Section~\ref{sec:eval:leak}).

\begin{table}[t]
  \caption{Verifiable ranked tallies. ``Public'' and ``talliers'' are what
  each party learns beyond the winner; ``paths'' are the witness paths
  of our certificate (Section~\ref{sec:security}).}
  \label{tab:comparison}
  \centering\scriptsize\setlength{\tabcolsep}{2.5pt}
  \begin{tabularx}{\columnwidth}{@{}lllll>{\raggedright\arraybackslash}X@{}}
    \toprule
    System & Public & Talliers & Proofs & Schulze & Machine-checked \\
    \midrule
    VSTTE'19~\cite{haines2019verifiable} & counts & counts & $\Sigma$ & yes & tally; crypto axiomatised \\
    VERICONDOR~\cite{vericondor2024} & counts & --- & $\Sigma$ & on counts & no \\
    Ordinos ext.~\cite{ordinos2020,ordinosext2021} & winner & nothing & MPC & yes & no \\
    Cortier et al.~\cite{cgy2022,yang2023thesis} & winner & nothing & MPC & yes, ties & UC proof of CGate \\
    Kryvos~\cite{kryvos2022} & winner & tally & SNARK & Condorcet & no \\
    this paper & winner, paths & counts & $\Sigma$ & yes, ties & ballot and certificate soundness, ZK \\
    \bottomrule
  \end{tabularx}
\end{table}
